Calibration needs a comparison available without labelled neighbours.
It uses the IDs returned by a larger budget on the same trained index.
Let $\mathcal R_\alpha(q)$ contain $k$ valid, unique returned IDs and
write $\mathcal R_{\rm ref}(q)=\mathcal R_{\alpha_{\rm ref}}(q)$. The
tested grid and reference are
\begin{equation}
 A=(200,100,64,32,16,8,4,2),\qquad \alpha_{\rm ref}=A[0]=200.
 \label{eq:6-grid}
\end{equation}
For the same $S$ queries $\mathcal Q_S$, count how many reference IDs
the smaller budget fails to return, ignoring result order and distances:
\begin{equation}
 \begin{aligned}
 D_S(\alpha)
 &=\sum_{q\in\mathcal Q_S}
   \left(k-|\mathcal R_\alpha(q)\cap\mathcal R_{\rm ref}(q)|\right)\\
 &=\frac12\sum_{q\in\mathcal Q_S}
   |\mathcal R_\alpha(q)\triangle\mathcal R_{\rm ref}(q)|.
 \end{aligned}
 \label{eq:6-agreement}
\end{equation}
The acceptance rule $D_S(\alpha)\leq\epsilon=1$ permits at most one
missing reference ID across the whole sample. The objective is to find
the smallest grid budget that passes:
\begin{equation}
 \widehat\alpha=\min\{\alpha\in A:D_S(\alpha)\leq\epsilon\}.
 \label{eq:6-selection-objective}
\end{equation}
If no smaller budget is accepted, the reference budget 200 is used.
\Cref{sec:6-3} states when bisection finds this minimum exactly.

Because the allowance stays at one ID, increasing the sample lowers
the permitted fraction of missing reference answers. For $k=10$,
\[
 \begin{array}{c|rrrr}S&32&64&128&500\\\hline
 \epsilon/(Sk)&0.003125&0.0015625&0.00078125&0.0002
 \end{array}
\]
The larger sample both covers more queries and imposes a stricter
relative threshold. The disjoint test queries in \cref{sec:7-5} measure
whether a budget chosen from the calibration pool retains quality.

\par\medskip\noindent\begin{minipage}{\linewidth}
\captionsetup{hypcap=false}
\centering
% Native schematic of the implemented decision and its offline evaluation.
\begin{tikzpicture}[
 font=\footnotesize\linespread{1}\selectfont,
 box/.style={draw=ThesisNavy!65,rounded corners=2pt,align=center,
 text width=2.8cm,minimum height=1.45cm,inner sep=4pt},
 arr/.style={-{Latex[length=2mm]},draw=ThesisNavy!80,thick}]
\node[box] (sample) at (0,0) {Calibration pool\\Even indices\\Sample $S$ queries};
\node[box,fill=ThesisNavy!7] (ref) at (3.6,0) {Reference search\\$\alpha_{\rm ref}=200$\\Top-$k$ IDs};
\node[box,fill=ThesisNavy!7] (test) at (7.2,0) {Bisection\\$D_S(\alpha)\leq1$\\Shared sample};
\node[box] (choice) at (10.8,0) {Chosen $\alpha$\\Caller applies it\\For batch reuse};
\node[box] (held) at (0,-2.0) {Held-out pool\\Odd query indices\\GT for evaluation};
\node[box,text width=6.4cm] (eval) at (7.2,-2.0)
 {Evaluate the selected arm\\Look up held-out recall and QPS\\from previously timed grid arms};
\draw[arr] (sample) -- (ref);
\draw[arr] (ref) -- (test);
\draw[arr] (test) -- (choice);
\draw[arr] (held) -- (eval);
\draw[arr,dashed] (choice.south) -- (10.8,-2.0) -- (eval.east);
\end{tikzpicture}

\captionof{figure}[Workload-level calibration and evaluation]{Workload-level calibration and evaluation.}
\label{fig:6-calibration}
\end{minipage}\par\medskip
Reference and trial searches in \cref{fig:6-calibration} use the same
$S$ queries. GT evaluates their choices afterwards on held-out queries;
selected-budget throughput comes from fixed-budget measurements.
\Cref{sec:a-budget-contract} defines the calls included in these timings.

\subsection{What agreement controls}

With nested survivors and fixed final ordering, disagreement has a
direct interpretation: it counts reference answers excluded by primary
screening. Under \Cref{prop:6-monotonicity},
\begin{equation}
 D_S(\alpha)=\sum_{q\in\mathcal Q_S}
       |\mathcal R_{\rm ref}(q)\setminus C_\alpha(q)|.
 \label{eq:6-reference-coverage}
\end{equation}
For any $k$-element outputs and GT set $\mathcal T_k$, the loss in GT
hits cannot exceed the missing reference IDs:
\[
 |\mathcal R_{\rm ref}\cap\mathcal T_k|-|\mathcal R_\alpha\cap\mathcal T_k|
 \leq|\mathcal R_{\rm ref}\setminus\mathcal R_\alpha|.
\]
Summing over the sample and dividing by $Sk$ yields
\begin{equation}
 \begin{aligned}
 r_{\mathcal Q_S}(\alpha)
 &\geq r_{\mathcal Q_S}({\rm ref})-\frac{D_S(\alpha)}{Sk},\\
 D_S(\alpha)\leq\epsilon
 &\implies r_{\mathcal Q_S}(\alpha)
       \geq r_{\mathcal Q_S}({\rm ref})-\frac{\epsilon}{Sk}.
 \end{aligned}
 \label{eq:6-relative-recall}
\end{equation}
This is a recall bound on the calibration sample relative to its
reference. It leaves the reference's own error to GT evaluation
(\cref{sec:7-recall-attribution}); it does not transfer automatically
to unseen queries.

\subsection{Sampling resolution and ordered testing}

A small sample can miss an infrequent change in returned IDs. The
following models describe how often such changes could remain unseen.
First count queries with any missing reference ID: call this $H_S$.
Since $H_S\leq D_S$, an accepted budget has $H_S\leq1$.
For a prespecified budget and independent, identically distributed
queries, let $u$ be that query-event probability. Its one-event
one-sided 95\% binomial limit satisfies \cite{nist_binomial}
\begin{equation}
 (1-u)^S+Su(1-u)^{S-1}=0.05.
 \label{eq:6-binomial-limit}
\end{equation}
A stronger model treats each returned-neighbour position as an
independent event with the same probability. Applying the limit to all
result slots gives the second row below:
\[
 \begin{array}{c|cc}
 &S=128&S=500\\\hline
 \text{i.i.d. query events}&0.03652&0.009452\\
 \text{i.i.d. homogeneous result slots}&0.003701&0.0009484
 \end{array}
\]
The slot model gives approximately $4.74/(Sk)$: at $S=500$ the scale is
about one missing reference result in a thousand; at $S=128$ it is about
four in a thousand. Queries being independent does not make their
result slots independent. These are model-based resolutions of reference
disagreement, rather than guarantees for the repeated fixed-pool samples.

Monotonicity allows budgets to be tested in a fixed sequence only when
each test has a valid statistical error bound \cite{learn_then_test}.
The implemented one-ID acceptance threshold checks agreement and does
not by itself certify future GT recall.

\paragraph{A finite-pool extension.}
For the actual sampling scheme, acceptance probability can instead be
computed from the composition of the calibration pool. Let
$X_q(\alpha)$ count the missing reference IDs of query $q$; assume it
stays fixed when that query moves between subsets or positions within
the execution regime. For pool size $N_{\rm cal}$, let
$n_j=|\{q:X_q=j\}|$. A uniformly sampled subset without replacement
passes only if all its queries have zero missing IDs, or exactly one
query has one missing ID \cite{nist_hypergeometric}:
\begin{equation}
 \Pr\{D_S(\alpha)\leq1\}
 =\frac{\binom{n_0}{S}+n_1\binom{n_0}{S-1}}{\binom{N_{\rm cal}}{S}}.
 \label{eq:6-finite-pool}
\end{equation}
Out-of-range binomial coefficients are zero. To turn acceptance
probabilities into GT-loss predictions also requires knowing which
budgets pass and their held-out quality. \Cref{sec:7-budget-trace}
checks the recorded subset sums against this calculation.
